\subsection{Formula}

Let $D$ be a district with area $A(D)$ and perimeter $P(D)$.
The Schwartzberg score is:
\begin{equation}
  S(D) = \frac{P(D)}{2\sqrt{\pi A(D)}}.
  \label{eq:schwartzberg}
\end{equation}
The denominator $2\sqrt{\pi A}$ is the perimeter of a circle with area $A$:
a circle of radius $r$ has area $\pi r^2$ and perimeter $2\pi r$,
so $2\pi r = 2\pi \sqrt{A/\pi} = 2\sqrt{\pi A}$.
Thus S is the ratio of the district's perimeter to the perimeter of a circle
with the same area.

\subsection{Properties}

\begin{proposition}
  $S(D) \geq 1$ for any polygon $D$, with equality if and only if $D$ is a disk.
\end{proposition}

\begin{proof}
  By the isoperimetric inequality, $P^2 \geq 4\pi A$, so $P \geq 2\sqrt{\pi A}$,
  giving $S = P/(2\sqrt{\pi A}) \geq 1$.
  Equality holds iff $P^2 = 4\pi A$, i.e., iff $D$ is a disk. \qed
\end{proof}

\begin{proposition}
  S is scale-invariant, rotation-invariant, and translation-invariant.
\end{proposition}

\begin{proof}
  Scaling by $\lambda$: $A \mapsto \lambda^2 A$ and $P \mapsto \lambda P$,
  so $S = \lambda P / (2\sqrt{\pi \lambda^2 A}) = P/(2\sqrt{\pi A})$. \qed
\end{proof}

\subsection{Canonical Values}

\begin{table}[ht]
\centering
\caption{Schwartzberg score for canonical shapes.}
\label{tab:s_canonical}
\begin{tabular}{lcc}
\toprule
Shape & $S$ & Approximate value \\
\midrule
Disk & $1$ & $1.000$ \\
Regular hexagon & $2 / (3^{3/4} \sqrt{\pi/3})$ & $1.050$ \\
Square of side $a$ & $2/\sqrt{\pi}$ & $1.128$ \\
Equilateral triangle & $2\sqrt[4]{3}/\sqrt{\pi}$ & $1.286$ \\
$1\times 3$ rectangle & $4\sqrt{3/\pi}/3$ & $1.302$ \\
$1\times 10$ rectangle & $11\sqrt{10/\pi}/10$ & $1.960$ \\
\bottomrule
\end{tabular}
\end{table}

Note: S $= 2/\sqrt{\pi} \approx 1.128$ for a square, derived from
$A = a^2$, $P = 4a$: $S = 4a / (2\sqrt{\pi a^2}) = 2/\sqrt{\pi}$.
